\section{What is a Nation}

\begin{quotex}
The existence of a nation is a daily plebiscite.
\flright{\textsc{Ernest Renan}}
\end{quotex}


In 1882, the one-time seminarian turned positivist historian, \textbf{Ernest Renan}, delivered his influential lecture ``What is a nation''. That the question even arises is a sign of modernity and Renan's ultimate answer seems far from the answer of Western ruling elite today. In the Ancient City, the question did not even arise. Renan confirms what both Fustel and Evola have claimed: the Ancient City was based on commonality of religion, race, and history cemented by the practice of daily rites, initiations and oaths of enduring fealty. Renan writes:

\begin{quotex} In the Ancient City religion was a matter of the very existence of the social group, which was an extension of the family... The religion of Athens was the cult of Athens itself, of its mythical founders, its laws and customs. It did not imply any dogmatic theology. This religion was a state religion. You were not an Athenian if you refused to practice it.
\end{quotex}


Nevertheless, old ideas die hard and some of the ancient ways persisted into the Middle Ages, though the protecting gods were replaced by the patron saint of the city or republic. The cult of that saint was important to the city and identified those loyal to it and those who were outsiders.

The idea of the divine origins of nations persisted in the Catholic Tradition. \textbf{Joseph de Maistre} denies that the nation is a rational construction. He explains:

\begin{quotex}
He who has not the power even of making an insect or a blade of grass, believes that he is the immediate author of sovereignty, the most important, sacred and fundamental thing in the moral and political world ... He believed that he could constitute nations ... that he could create national unity ... It is a truth as certain in its way as a mathematical proposition that no great institution results from deliberation, and that human works are fragile in proportion to the number of men concerned in them and the degree to which science and reasoning have been used \emph{a priori}.
\end{quotex}


\textbf{St Pope Pius X} asserts: ``The Creator of mankind is also the Founder of human societies.'' A nation is a living organism, capable of growth and activity, decline and decay. They also have vocations, as Pius XII wrote:

\begin{quotex}
Like individuals, peoples are prosperous or unfortunate, influential or obscurely sterile to the degree that they are docile or rebellious to their vocations. ... . \emph{There was no surer sign of moral disorder than a decline in love for the national fatherland}.
\end{quotex}


Without going into all the interesting historical details of Renan's lecture, he analyzes several factors that could define nationhood in modern times, only to reject them all. These are: ``race, language, interests, religious affinity, geography, military necessity.'' Renan comes to the conclusions:

\begin{quotex}
A nation is a spiritual principle resulting from the profound complications of history, \emph{a spiritual family and not a group determined by the configuration of the soil}.
\end{quotex}


Renan concludes the lecture with this understanding of the inner nature of the nation:

\begin{quotationx}
A nation is a soul, a spiritual principle. Two things that, in truth, are only one, make up this soul, this spiritual principle. One of these lies in the past \textbf{Continuity} , the other in the present \textbf{Solidarity} . One is the possession in common of a rich legacy of memories; the other is the present consent, the desire to live together, the will to continue to validate the heritage that has been jointly received. Man is not improvised. The nation, like the individual, is the culmination for a long past of effort, sacrifice and devotion. This makes the cult of ancestors all the more legitimate; it is our ancestors who made us what we are. A heroic past, great men, glory (genuine glory, I mean), that is the social capital on which a national idea is founded. To have common glories in the past, and a common will in the present; to have done great things together, and to seek to do so again, those are the essential conditions for being a people. One loves in proportion to the sacrifices to which one has consented, the evils that one has suffered. One loves the house that one has built and passes on.

A nation is thus a great solidarity ... It presupposes a past, yet it is summed up in the present by a tangible fact: \emph{the clearly expressed consent and desire to continue a common life}.
\end{quotationx}



\flrightit{Posted on 2011-07-26 by Cologero}

